\bookchapter{الفصل 10: الدفع عبر الهاتف المحمول والخدمات الحياتية}

\chapterintro{بصراحة، كان جيل والديّ في البداية متردداً تجاه فكرة “الدفع بالهاتف”. “إذا كانت الأموال في الهاتف، ماذا لو سرقه أحد؟” هذا القلق طبيعي، لكن الواقع هو: احتمال ضياع النقد في جيبك أعلى بكثير من احتمال اختراق الدفع عبر الهاتف. الدفع عبر الهاتف لا يعني وضع أموالك في الهاتف، بل هو مجرد مفتاح — مفتاح يفتح حسابك البنكي.}

\booksection{أولاً: كيف تربط البطاقة البنكية فعلياً؟}

سواء كنت تستخدم Alipay أو WeChat، فإن ربط البطاقة لأول مرة يتطلب هذه الخطوات:

\begin{enumerate}
  \item افتح التطبيق، وابحث عن "حسابي" أو "المحفظة". انقر على "إضافة بطاقة بنكية".
  \item أدخل رقم البطاقة. تدعم معظم التطبيقات الآن التعرف على رقم البطاقة بالتصوير المباشر، دون الحاجة إلى إدخال الأرقام يدوياً.
  \item سيرسل البنك رسالة نصية برمز تحقق إلى رقم هاتفك المسجل. أدخل الرمز، واضبط رمز دفع مكوّن من 6 أرقام.
  \item من الأفضل أن يكون رمز الدفع هذا مختلفاً عن رمز بطاقتك البنكية ورمز قفل الشاشة — ثلاثة رموز مختلفة، ولا تتكاسل وتجعلها متطابقة.
\end{enumerate}

تلميح صغير: إذا كنت قلقاً بشأن الأمان، يمكنك فتح حساب إلكتروني من الفئة الثانية أو الثالثة في تطبيق البنك، مخصصاً لربط الدفع عبر الهاتف. ضع في هذه البطاقة مصاريفك اليومية فقط، فحتى لو حدث خطأ، ستكون الخسارة محدودة.

\booksection{ثانياً: رمز الدفع هو “من يراه يستلم المال”}

كثير من الناس يخلطون بين رمز الدفع ورمز الاستلام، وهذا الفرق مهم جداً.

\begin{itemize}
  \item \textbf{رمز الدفع}: الرمز الشريطي + رمز QR الذي تعرضه على أمين الصندوق لمسحه. هذا الرمز يمثل "محفظتك الإلكترونية"، فمن يمسحه يستطيع سحب الأموال من حسابك. لذا\textbf{لا ترسل لقطة شاشة منه إلى أي شخص أبداً}. يتجدد هذا الرمز تلقائياً كل دقيقة، لمنع استخدامه من قبل الآخرين.
  \item \textbf{رمز الاستلام}: عندما يمسحه الآخرون، يحوّلون الأموال إليك. يمكن أخذ لقطة شاشة لهذا الرمز وحفظها وإرسالها للآخرين، وهو آمن.
\end{itemize}

تعرّف على هذا الفرق، فهو يقيك من الكثير من فخاخ الاحتيال. كل من يطلب منك “إرسال لقطة شاشة لرمز الدفع” هو محتال بنسبة 100\%.

\booksection{ثالثاً: هذه الإعدادات الأمنية، 5 دقائق لتضبطها توفر عليك الكثير من القلق}

\begin{itemize}
  \item \textbf{الدفع بالبصمة/الوجه}: في إعدادات الدفع في Alipay وWeChat، فعّل الدفع البيومتري. هكذا يتطلب كل دفع ضغط البصمة أو مسح الوجه، فحتى لو فقدت هاتفك، لن يتمكن من يجده من إنفاق أموالك.
  \item \textbf{قفل الأمان}: لدى Alipay وWeChat وظيفة "قفل الأمان". عند تفعيلها، يتطلب الدخول إلى صفحة "المحفظة" في التطبيق تحقق إضافي من الهوية. المسار: Alipay حسابي الإعدادات الحساب والأمان مركز الأمان قفل الأمان.
  \item \textbf{فحص الدفع بدون رمز}: كثيرون اشتركوا في خدمات دفع بدون رمز ونسوها. اذهب إلى "إعدادات الدفع" في Alipay وWeChat، وابحث عن "الدفع بدون رمز/الخصم التلقائي"، لترى الخدمات المشترك فيها، وألغِ ما لا تحتاجه.
  \item \textbf{إدارة أجهزة تسجيل الدخول}: راجع دورياً الأجهزة التي سجّلت الدخول إلى حسابك. إذا اكتشفت جهازاً غير معروف، اطرده فوراً. المسار: Alipay الإعدادات الحساب والأمان إدارة أجهزة تسجيل الدخول.
\end{itemize}

\booksection{رابعاً: الحافلات والمترو: الهاتف هو بطاقة المواصلات}

قلة قليلة الآن تستخدم بطاقة مواصلات مادية عند ركوب الحافلات والمترو. لتمرير الهاتف على بوابة الدفع، هناك طريقان:

\begin{itemize}
  \item \textbf{بطاقة مواصلات افتراضية في محفظة الهاتف}: افتح بطاقة مواصلات افتراضية مباشرة في Apple Wallet أو Huawei Wallet أو Xiaomi Wallet. ميزة هذا الحل: عند المرور لا تحتاج إلى فتح قفل الهاتف، بل قرّبه فحسب، وسرعة الاستجابة هي الأسرع، ويمكن استخدامه دون اتصال بالإنترنت. بطاقات المواصلات الافتراضية في بعض المدن تتمتع بنفس خصومات التحويلات للبطاقات المادية.
  \item \textbf{رمز ركوب Alipay/WeChat}: ابحث في التطبيق عن "رمز الركوب"، وفعّل رمز ركوب الحافلات أو المترو في مدينتك. ميزة هذا الحل: لا تحتاج إلى شحن، إذ يُخصم المبلغ مباشرة من الرصيد أو البطاقة البنكية المرتبطة. لكنه يتطلب اتصالاً بالإنترنت، وعليك فتح التطبيق للعثور على الرمز قبل التمرير.
\end{itemize}

أنا شخصياً أفضل بطاقة المواصلات الافتراضية في محفظة الهاتف. رغم أنك تحتاج إلى شحن 10 أو 20 يواناً في المرة الأولى، إلا أنها أسهل بكثير في الاستخدام، خاصة عند اللحاق بالمترو، فلا تحتاج إلى التوقف لفتح قفل الهاتف والبحث عن التطبيق.

\booksection{خامساً: فواتير الخدمات، توصيل الطعام، وطلب السيارات}

\begin{itemize}
  \item \textbf{دفع فواتير الماء والكهرباء والغاز}: في "دفع الفواتير" في Alipay أو WeChat، اربط أرقام حسابات الماء والكهرباء والغاز في منزلك. وعند استلام الفاتورة كل شهر، انقر عليها مباشرة للدفع. يمكن العثور على رقم الحساب عادةً في الفاتورة الورقية، أو بالاتصال بخدمة العملاء.
  \item \textbf{مقارنة أسعار توصيل الطعام}: الأسعار والعروض في المتجر نفسه على Meituan وEle.me تختلف كثيراً. انظر في كليهما قبل الطلب، فقد توفر بضعة يوانات أحياناً. كما أن وظيفة "الطلب الجماعي" في Meituan، عند الطلب مع الأصدقاء، تقلّل السعر كثيراً.
  \item \textbf{أمان طلب السيارات}: في Didi أو Amap Ride-hailing، تذكر ضبط "جهة الاتصال للطوارئ" و"مشاركة الرحلة". عند طلب السيارة، يشارك التطبيق رحلتك تلقائياً مع جهة الاتصال التي حددتها. كما يُفضّل في الطلبات الليلية اختيار "سيارة أجرة" بدلاً من "سيارة سريعة"، فسيارات الأجرة تخضع لرقابة صناعية، ومستوى الأمان أعلى.
\end{itemize}

\booksection{سادساً: التأمين الصحي والضمان الاجتماعي والخدمات الحكومية عبر الهاتف}

التغيّرات في السنوات الأخيرة كبيرة، فالكثير من الأمور التي كانت تتطلب الذهاب إلى الصالة والوقوف في طابور، يمكن إنجازها الآن في دقائق عبر الهاتف.

\begin{itemize}
  \item \textbf{الشهادة الإلكترونية للتأمين الصحي}: ابحث في Alipay أو WeChat عن "الشهادة الإلكترونية للتأمين الصحي"، وبعد تفعيلها، يمكنك عند زيارة الطبيب أو شراء الدواء من الصيدلية استخدام الهاتف مباشرة، دون حمل البطاقة المادية. يتطلب أول استخدام التحقق بتعرف الوجه، ثم يُنشأ رمز ديناميكي عند كل استخدام.
  \item \textbf{الاستعلام عن الضمان الاجتماعي}: نزّل تطبيق "Palm 12333" أو استخدم "مركز المواطن" في Alipay، للاستعلام عن سجل مدفوعات الضمان الاجتماعي ورصيد حساب المعاش والمطالبة ببدل البطالة وغير ذلك.
  \item \textbf{ضريبة الدخل}: التسوية السنوية لضريبة الدخل من مارس إلى يونيو يمكن إنجازها على تطبيق "ضريبة الدخل الفردي". تذكر التحقق من الخصومات الإضافية الخاصة التي يمكن الإقرار عنها (تعليم الأبناء، فوائد قروض الإسكان، إعالة كبار السن وغيرها)، فكلها تساعدك على تقليل الضريبة أو استردادها.
\end{itemize}

\begin{tcbwarn}{تذكر هذه القواعد، لتقي نفسك من الاحتيال}
\begin{itemize}
  \item الشرطة والنيابة والبنوك وخدمة العملاء لن يطلبوا منك تحويل الأموال إلى "حساب آمن".
  \item أي شخص يطلب منك إرسال لقطة شاشة لرمز الدفع هو محتال.
  \item لا تنقر على الروابط غير المعروفة، خاصة الرسائل النصية من نوع "انتهاء صلاحية ETC" أو "إيقاف بطاقة الضمان الاجتماعي".
  \item المبالغ المستردة من التسوق عبر الإنترنت تُعاد عبر المسار الأصلي، ولا تحتاج إلى تقديم رقم بطاقتك البنكية.
\end{itemize}

\end{tcbwarn}

\begin{tcbpractice}{افعل ذلك الآن: أجرِ فحصاً أمنياً لتطبيقات الدفع لديك}
\begin{enumerate}
  \item افتح Alipay وWeChat، وتأكد من تفعيل الدفع بالبصمة/الوجه.
  \item ابحث عن وظيفة "قفل الأمان" وفعّلها.
  \item راجع قائمة "الدفع بدون رمز/الخصم التلقائي"، وألغِ الخدمات التي لم تعد بحاجة إليها.
  \item اعرض قائمة "أجهزة تسجيل الدخول"، وأزِل الأجهزة غير المعروفة.
  \item فعّل شهادتك الإلكترونية للتأمين الصحي، حتى لا تضطر لحمل البطاقة المادية في زيارتك القادمة للمستشفى.
\end{enumerate}

\end{tcbpractice}
